\subsubsection{Vectors and three-dimensional geometry}
For vectors $\vec a,\vec b\in\mathbb R^3$,
\begin{align}
\vec a\cdot\vec b&=a_1b_1+a_2b_2+a_3b_3,\quad \vec a\perp\vec b\iff \vec a\cdot\vec b=0. \tag{1.1.1.1}\\
\vec a\times\vec b&=\begin{pmatrix}a_2b_3-a_3b_2\\a_3b_1-a_1b_3\\a_1b_2-a_2b_1\end{pmatrix}. \tag{1.1.1.2}\\
\text{unit vector along }\vec v&=\frac{\vec v}{\norm{\vec v}},\qquad \norm{c\vec v}=|c|\norm{\vec v}. \tag{1.1.1.3}
\end{align}
The vector product is normal to both vectors in the product. Distances and line/plane relations are
\begin{align}
d(\text{parallel lines})&=\frac{\norm{(\vec b-\vec a)\times\vec v}}{\norm{\vec v}}. \tag{1.1.2.1}\\
d(\text{skew lines})&=\frac{|(\vec b-\vec a)\cdot(\vec v_1\times\vec v_2)|}{\norm{\vec v_1\times\vec v_2}}. \tag{1.1.2.2}\\
\text{line parallel to plane}&\iff \vec d\cdot\vec n=0. \tag{1.1.2.3}\\
\text{plane through }\vec r_0\text{ with normal }\vec n:\quad \vec n\cdot(\vec r-\vec r_0)&=0. \tag{1.1.2.4}\\
\text{line with point }\vec A\text{ and direction }\vec d:\quad \vec r&=\vec A+t\vec d. \tag{1.1.2.5}
\end{align}

\subsubsection{Matrices and determinants}
\begin{align}
\det(AB)&=\det(A)\det(B),\qquad \det(A^n)=(\det A)^n. \tag{1.2.1.1}\\
A^{-1}&=\frac{\operatorname{adj}(A)}{\det A},\qquad \det A\ne0. \tag{1.2.1.2}\\
A\vec x=\vec b&\quad\Longrightarrow\quad \vec x=A^{-1}\vec b\quad(A\text{ invertible}). \tag{1.2.1.3}\\
[A\mid \vec b]&\xrightarrow{\text{elementary row operations}}[I\mid\vec x]\quad\text{(Gaussian elimination)}. \tag{1.2.1.4}
\end{align}
The identity used to factor the determinant in Q36 is
\[
a^2+b^2+c^2-ab-bc-ca=\frac12\bigl((a-b)^2+(b-c)^2+(c-a)^2\bigr).\tag{1.2.2.1}
\]

\subsubsection{Differentiation and differential equations}
\begin{align}
\frac{dy}{dx}&=\frac{dy/dt}{dx/dt},\qquad \frac{d^2y}{dx^2}=\frac{d}{dt}(dy/dx)\bigg/\frac{dx}{dt}. \tag{1.3.1.1}\\
 y'+P(x)y&=Q(x),\quad \mu(x)=e^{\int P(x)\,dx},\quad \frac{d}{dx}(\mu y)=\mu Q. \tag{1.3.1.2}\\
y=vx&\quad\Longrightarrow\quad y'=v+xv'\quad\text{(homogeneous first-order ODE)}. \tag{1.3.1.3}\\
(1+(y')^2)^3&=r^2(y'')^2\quad\text{for a circle family of fixed radius }r. \tag{1.3.1.4}\\
\frac{d}{dx}\log|u(x)|&=\frac{u'(x)}{u(x)},\qquad u(x)\ne0. \tag{1.3.1.5}
\end{align}

\subsubsection{Integration and area}
\begin{align}
\int\frac{du}{a^2-u^2}&=\frac1{2a}\log\left|\frac{a+u}{a-u}\right|+C. \tag{1.4.1.1}\\
\int\frac{du}{1+u^2}&=\tan^{-1}u+C. \tag{1.4.1.2}\\
\int\frac{u\,du}{\sqrt{1+u^2}}&=\sqrt{1+u^2}+C. \tag{1.4.1.3}\\
A&=\int_a^b |y_{\mathrm{upper}}-y_{\mathrm{lower}}|\,dx. \tag{1.4.1.4}\\
\frac{P(x)}{(x-a)(x^2+b^2)}&=Q(x)+\frac{A}{x-a}+\frac{Bx+C}{x^2+b^2}\quad\text{(partial fractions after division)}. \tag{1.4.1.5}\\
\int u\,dv&=uv-\int v\,du. \tag{1.4.1.6}
\end{align}

\subsubsection{Functions, relations and operations}
At an interior point, continuity requires matching one-sided limits; differentiability additionally requires matching one-sided derivatives:
\[\lim_{x\to a^-}f(x)=f(a)=\lim_{x\to a^+}f(x),\qquad f'_{-}(a)=f'_{+}(a)=f'(a).\tag{1.5.1.1}\]
\[
\begin{aligned}
\text{reflexive: }&(a,a)\in R,\\
\text{symmetric: }&(a,b)\in R\Rightarrow(b,a)\in R,\\
\text{transitive: }&(a,b),(b,c)\in R\Rightarrow(a,c)\in R.
\end{aligned}\tag{1.5.1.2}
\]
\[\text{For componentwise addition on }\mathbb R^2:\quad e=(0,0),\quad (a,b)^{-1}=(-a,-b).\tag{1.5.1.3}\]
\[a*b=(a+b)\bmod n,\quad 0\text{ is the identity},\quad (-a)\bmod n\text{ is the additive inverse}.\tag{1.5.1.4}\]

\subsubsection{Probability}
\begin{align}
P(A\mid E)&=\frac{P(A\cap E)}{P(E)}. \tag{1.6.1.1}\\
P(A\mid E)&=\frac{P(E\mid A)P(A)}{P(E)}. \tag{1.6.1.2}\\
X\sim\operatorname{Bin}(n,p):\quad P(X=k)&=\binom nkp^k(1-p)^{n-k}. \tag{1.6.1.3}\\
G(z)&=\sum_jP(D=j)z^j. \tag{1.6.1.4}\\
G_5(z)&=(0.4z+0.6z^{-1})^5 \text{ for the five-step walk}. \notag\\
E[X]&=\sum_x xP(X=x),\quad \operatorname{Var}(X)=E[X^2]-E[X]^2. \tag{1.6.1.5}\\
P(X=k)&=\frac{\binom{K}{k}\binom{N-K}{n-k}}{\binom Nn}. \tag{1.6.1.6}\\
P(E)&=\sum_iP(E\mid A_i)P(A_i). \\
&\text{(law of total probability).} \tag{1.6.1.7}
\end{align}

\subsubsection{Linear programming}
\[\text{Evaluate the objective at each feasible corner; an optimum occurs at a feasible vertex.}\tag{1.7.1.1}\]

\subsubsection{Coordinate geometry and optimisation}
The tangent to $x^2+y^2=r^2$ at $(r\cos\theta,r\sin\theta)$ is
\[
x\cos\theta+y\sin\theta=r.\tag{1.8.1.1}
\]
For the ellipse $x^2/a^2+y^2/b^2=1$, a parametrisation is $(a\cos t,b\sin t)$, and triangle area is one half base times perpendicular height:
\[\text{ellipse point }(a\cos t,b\sin t),\quad A_{\triangle}=\frac12 bh.\tag{1.8.1.2}\]

\subsubsection{Engineering methods: choosing the right tool}
Matrix methods are used for vector and matrix conditions and for Gaussian elimination. The five-step random walk and the small binomial problem use probability-generating polynomial coefficients where it makes the probability structure clearer. Laplace transforms are useful for suitable linear time-domain differential equations with initial conditions; they are not forced into the nonlinear or variable-coefficient differential equations in this set where they would not simplify the solution.
